% ─────────────────────────────────────────────────────────────────
% CENG 467 Midterm – Question 1: Text Classification
% Student ID: 310201050
% ─────────────────────────────────────────────────────────────────

\section{Question 1: Representation Learning in Text Classification}

\subsection{Dataset}

The IMDb Large Movie Review Dataset \cite{maas2011imdb} consists of 50{,}000
binary-labelled movie reviews (25{,}000 train / 25{,}000 test), evenly split
between positive and negative sentiment classes. Because fine-tuning a
Transformer on the full training corpus exceeds the available memory budget
of the Apple~M2 (8\,GB RAM), a reproducible subset was drawn using seed 3102
derived from the author's student identifier: $6{,}000$ training samples,
$1{,}000$ validation samples, and $1{,}000$ held-out test samples. The test
set was used \emph{exactly once} for final evaluation.

\subsection{Preprocessing and Tokenisation Strategies}

A shared normalisation pipeline is applied to all texts before
model-specific tokenisation. HTML tags are stripped with a regular
expression, the text is lower-cased, all non-alphabetic characters are
removed, and consecutive whitespace is collapsed to a single space. Three
tokenisation strategies are then applied, each suited to the
representational assumptions of a different model.

\paragraph{Strategy A — Whitespace Split with Stopword Removal.}
Tokens are produced by splitting on whitespace after normalisation,
followed by removal of English stopwords (NLTK list, 179 tokens). This
yields compact, content-bearing token sequences and is used by the TF-IDF
representation. Removing stopwords reduces feature noise but discards
negation words such as \textit{not}; the TF-IDF bigram range partially
compensates for this loss.

\paragraph{Strategy B — NLTK \texttt{word\_tokenize}.}
NLTK's rule-based tokeniser is applied after normalisation \emph{without}
stopword removal, preserving punctuation-adjacent morphology and
contractions that carry affective weight (e.g.\ \textit{wasn't},
\textit{couldn't}). This strategy is used by BiLSTM to retain richer
sequential cues.

\paragraph{Strategy C — WordPiece subword tokenisation.}
DistilBERT employs its own WordPiece subword tokeniser via
\texttt{DistilBertTokenizerFast}, which decomposes out-of-vocabulary words
into known sub-units, making it the most robust of the three approaches.

\subsection{Models}

\subsubsection{TF-IDF + Logistic Regression}

A sparse TF-IDF matrix is constructed from Strategy-A tokens with $30{,}000$
maximum features, unigram and bigram ranges, and sublinear TF scaling
($\text{tf} \leftarrow 1+\log\text{tf}$). A Logistic Regression classifier
with $C=0.7$ is trained via L-BFGS. This constitutes the \emph{sparse}
representation baseline.

\subsubsection{BiLSTM}

A two-layer bidirectional LSTM \cite{hochreiter1997lstm} is implemented with
the configuration shown in Table~\ref{tab:bilstm_hp}. The final hidden
states from the last forward and backward layers are concatenated and
passed through a linear classification head. This constitutes the
\emph{dense sequential} representation.

\begin{table}[h]
\centering
\caption{BiLSTM Hyperparameters}
\label{tab:bilstm_hp}
\begin{tabular}{ll}
\toprule
Hyperparameter & Value \\
\midrule
Vocabulary size     & 20{,}000 \\
Embedding dimension & 128 \\
Hidden dimension    & 128 \\
Number of layers    & 2 (bidirectional) \\
Dropout             & 0.35 \\
Maximum sequence length & 256 \\
Batch size          & 32 \\
Epochs              & 5 \\
Optimiser / LR      & Adam / $10^{-3}$ \\
Gradient clipping   & 1.0 \\
\bottomrule
\end{tabular}
\end{table}

\subsubsection{DistilBERT}

DistilBERT-base-uncased \cite{sanh2019distilbert} is fine-tuned for sequence
classification (Table~\ref{tab:bert_hp}). The full self-attention stack
provides \emph{contextual} representations where each token's embedding is
conditioned on the entire input.

\begin{table}[h]
\centering
\caption{DistilBERT Fine-Tuning Hyperparameters}
\label{tab:bert_hp}
\begin{tabular}{ll}
\toprule
Hyperparameter & Value \\
\midrule
Pretrained model        & \texttt{distilbert-base-uncased} \\
Maximum sequence length & 256 \\
Batch size              & 16 \\
Epochs                  & 3 \\
Learning rate           & $2 \times 10^{-5}$ \\
Weight decay            & 0.01 \\
Warmup ratio            & 10\% \\
Gradient clipping       & 1.0 \\
Scheduler               & Linear with warmup \\
\bottomrule
\end{tabular}
\end{table}

\subsection{Evaluation Results}

All models are evaluated on the same held-out test split ($n=1{,}000$) using
Accuracy and Macro-F1. Given the balanced class distribution, both metrics
yield consistent estimates of generalisation
(Table~\ref{tab:q1_results}).

\begin{table}[h]
\centering
\caption{Test-Set Performance on IMDb (seed 3102)}
\label{tab:q1_results}
\begin{tabular}{lcc}
\toprule
Model & Accuracy & Macro-F1 \\
\midrule
TF-IDF + Logistic Regression & 0.8550 & 0.8548 \\
BiLSTM                       & 0.7490 & 0.7488 \\
DistilBERT                   & \textbf{0.8970} & \textbf{0.8970} \\
\bottomrule
\end{tabular}
\end{table}

The pattern is striking: the sparse TF-IDF baseline (85.5\%) substantially
outperforms the dense BiLSTM (74.9\%), while the contextual DistilBERT
(89.7\%) outperforms both. The strong showing of TF-IDF reflects the
fact that IMDb reviews are long and lexically rich in sentiment-bearing
unigrams, exactly the regime in which sparse linear models excel. The
BiLSTM, despite its sequential modelling capacity, is data-hungry: with
only 6{,}000 randomly initialised embeddings, it cannot generalise as
well as a well-regularised linear model on the same data.

\subsection{Misclassification Analysis}

Five misclassified examples produced by DistilBERT were inspected in
detail; the full text of each is preserved in
\texttt{q1\_results.json}. Three of the five errors share a single
recurring pattern.

\paragraph{Pattern 1 — Ironic praise of niche or ``cult'' material.}
Reviews [8], [14], and [31] are all positive reviews predicted as
negative. They share a rhetorical strategy in which the reviewer
enthusiastically endorses a film while using ironic or self-deprecating
descriptors (\textit{``Citizen Cane this film is not''}, \textit{``a
mess''}, \textit{``bad hair''}, \textit{``about as good as a bad
dino-flick can get''}). DistilBERT picks up on the surface-level negative
adjectives and misclassifies the overall stance, which is unmistakably
positive on a careful reading.

\paragraph{Pattern 2 — Mixed sentiment with a reversed conclusion.}
Review [50] is a negative review predicted as positive: the writer opens
by criticising the film as boring and predictable, then spends most of
the body praising specific actors and on-screen chemistry, before
returning to the negative verdict. DistilBERT's 256-token window captures
predominantly the central praise, missing the bracketing critical frame.

\paragraph{Pattern 3 — Affective negation with mixed lexical signal.}
Review [1] explicitly disclaims itself (\textit{``this film will never
win an oscar''}, \textit{``Citizen Cane this film is not''}) but
expresses genuine enjoyment (\textit{``I had fun''}, \textit{``I liked
the cheese''}). The lexical signal is genuinely conflicted, and the
model resolves it incorrectly.

These three patterns —- irony, mixed-sentiment with reversal, and
affective negation —- are all variants of a single underlying problem:
\emph{the model is reading lexical polarity rather than rhetorical
stance}. Such errors are well-documented for fine-tuned Transformers and
are not eliminable without explicit modelling of discourse structure or
substantially larger training data.

\subsection{Discussion: Representation Types}

\paragraph{TF-IDF (Sparse).} Despite its simplicity, TF-IDF with logistic
regression delivered $85.5\%$ accuracy --- well within ten points of
the DistilBERT result --- because IMDb reviews are long and
sentiment-rich at the unigram level. The model is computationally cheap,
fully interpretable through feature weights, and requires no GPU.
However, it is completely insensitive to word order and context; surface
forms such as \textit{good} and \textit{not good} receive independent
representations.

\paragraph{BiLSTM (Dense Sequential).} Surprisingly, the BiLSTM
\emph{underperforms} TF-IDF on this subset by more than 10 points. This
counterintuitive ordering reflects two factors. First, the BiLSTM is
initialised with random embeddings, so it must learn its representations
from only 6{,}000 examples --- insufficient to outperform a
well-regularised sparse linear model. Second, the long-tail vocabulary
of IMDb reviews is partially clipped to 20{,}000 types, and rare but
sentiment-bearing words are mapped to \texttt{<UNK>}. With pre-trained
GloVe or word2vec initialisation the result would likely improve, but
even then the gap to DistilBERT would persist.

\paragraph{DistilBERT (Contextual).} Pre-trained self-attention provides
rich, context-dependent representations that markedly outperform both
sparse and recurrent baselines. Each token is encoded with global
context at no additional training cost beyond fine-tuning. The
trade-off is computational: DistilBERT requires substantially more
memory and training time than the other two approaches. On the M2 the
3-epoch fine-tune took roughly $45$\,minutes, compared with seconds for
the sparse model.